\section{Datengrundlage und Toolkontext}
\label{sec:results_context}

Dieses Unterkapitel beschreibt die Datengrundlage der Untersuchung sowie den technischen Kontext der durchgeführten \textit{VibeCoding-Challenge}. Die Analyse basiert auf einem Datensatz von $N=18$ Probanden, die im Rahmen einer kontrollierten Coding-Challenge Aufgaben zur Entwicklung einer Todo-Applikation bearbeiteten.

\subsection{Studiensetting und Rahmenbedingungen}
Die Studie wurde unter einheitlichen Rahmenbedingungen durchgeführt, um die Vergleichbarkeit der Ergebnisse zu gewährleisten:
\begin{itemize}
    \item \textbf{Zeitlimit:} Alle Probanden hatten ein striktes Zeitlimit von 60 Minuten für die Bearbeitung der Aufgaben.
    \item \textbf{Tooling:} Als Entwicklungsumgebung diente Visual Studio Code mit der GitHub Copilot Extension.
    \item \textbf{KI-Modell:} Als zugrundeliegendes Large Language Model (LLM) wurde \textbf{Claude 3.5 Sonnet} verwendet. Die Interaktion erfolgte über das Chat-Interface von GitHub Copilot, wobei das Modell explizit als Backend konfiguriert war.
    \item \textbf{Aufgabenstellung:} Die Probanden erhielten eine Liste von 8 Aufgaben (T1--T8) mit steigender Komplexität, beginnend bei Backend-Persistenz bis hin zu komplexen Frontend-Features und Export-Funktionen.
\end{itemize}

\subsection{Datenquellen und Artefakte}
Die empirische Auswertung stützt sich auf drei primäre Datenquellen, die eine Triangulation der Ergebnisse ermöglichen:

\begin{enumerate}
    \item \textbf{Code-Artefakte (Git-Repositories):}
    Die finalen Implementierungen der Probanden liegen als separate Git-Branches vor. Diese Branches enthalten den vollständigen Quellcode zum Zeitpunkt des Zeitablaufs.
    \begin{itemize}
        \item Pfad: Remote-Branches mit dem Suffix \texttt{*-vibe} (z.B. \texttt{origin/august-martin-vibe}).
    \end{itemize}

    \item \textbf{Chat-Logs:}
    Die vollständigen Interaktionsprotokolle zwischen Proband und KI-Assistent wurden exportiert. Diese Protokolle enthalten alle Prompts der Nutzer sowie die Antworten und Code-Generierungen der KI.
    \begin{itemize}
        \item Pfad: \texttt{thesis/chats/} (Format: \texttt{.docx}).
    \end{itemize}

    \item \textbf{Survey-Daten:}
    Im Anschluss an die Challenge füllten die Probanden einen Fragebogen zur Selbsteinschätzung und subjektiven Wahrnehmung der KI-Interaktion aus.
    \begin{itemize}
        \item Pfad: \texttt{thesis/surveys/Probanden-Einstufungsformular (Responses).xlsx}.
    \end{itemize}
\end{enumerate}

\subsection{Dateninventar und Vollständigkeit}
Tabelle \ref{tab:data_inventory} zeigt das Inventar der verfügbaren Datenartefakte. Insgesamt liegen für 18 Probanden Daten vor.

\begin{table}[h]
    \centering
    \begin{tabular}{|l|c|l|}
    \hline
    \textbf{Artefakt-Typ} & \textbf{Anzahl} & \textbf{Anmerkung} \\
    \hline
    Git-Branches (\texttt{*-vibe}) & 16 & 2 Probanden arbeiteten lokal/ohne Push* \\
    Chat-Logs & 18 & Vollständig für alle Teilnehmer \\
    Survey-Einträge & 18 & Vollständig ($N=18$) \\
    \hline
    \end{tabular}
    \caption{Inventar der Datenartefakte. (*Für die fehlenden Remote-Branches liegen lokale Stände oder Chat-Rekonstruktionen vor, die für die Analyse genutzt wurden.)}
    \label{tab:data_inventory}
\end{table}

\subsection{Zuordnung und Zusammenführung der Daten}
Die Zusammenführung der Datenquellen erfolgte über eine Mapping-Tabelle, die die Klarnamen aus dem Survey und den Chat-Dateinamen den entsprechenden Git-Branches zuordnet.
\begin{itemize}
    \item \textbf{Survey $\leftrightarrow$ Branch:} Das Survey-Feld "Branch-Name" sowie der Name des Probanden dienten als Schlüssel.
    \item \textbf{Chat $\leftrightarrow$ Branch:} Die Dateinamen der Chat-Logs (z.B. \texttt{Max\_Helling.docx}) wurden normalisiert und den entsprechenden Branches (z.B. \texttt{helling-max-vibe}) zugeordnet.
\end{itemize}
Die Survey-Daten enthalten zudem eine Einteilung in Kompetenzgruppen (No-Code, Low-Code, Mid-Code, High-Code), die in der späteren Analyse als vergleichendes Merkmal herangezogen wird.

\subsection{Reproduzierbarkeit und Auswertungswerkzeuge}
Um die Ergebnisse objektiv und reproduzierbar zu gestalten, kommen automatisierte Analysewerkzeuge zum Einsatz, deren Resultate im Ordner \texttt{evaluation/} persistiert werden:

\begin{itemize}
    \item \textbf{Statische Code-Analyse:} Einsatz von \texttt{ESLint} zur Überprüfung der Code-Qualität und Einhaltung von Standards.
    \item \textbf{Automatisierte Tests:} Ausführung der vorhandenen Test-Suite (basierend auf \texttt{Jest/Vitest}), um die funktionale Korrektheit der Implementierungen (Task Completion) zu verifizieren.
    \item \textbf{Chat-Parsing:} Algorithmische Auswertung der Chat-Logs zur Bestimmung von "Attempted"-Tasks und Identifikation von Prompting-Mustern.
\end{itemize}

\subsection{Abgrenzung zu Kapitel 7}
Das vorliegende Kapitel 6 beschränkt sich auf die deskriptive Darstellung der Ergebnisse und die objektive Auswertung der Daten. Eine Interpretation dieser Ergebnisse, die Diskussion von Ursache-Wirkungs-Beziehungen sowie die Beantwortung der Forschungsfragen erfolgen im anschließenden Kapitel 7 (Diskussion).
